\documentclass[11pt]{article}
\usepackage[a4paper,margin=1in]{geometry}
\usepackage{amsmath,amssymb}
\usepackage[T1]{fontenc}
\usepackage{lmodern}
\usepackage{microtype}
\usepackage{xcolor}

\setlength{\parindent}{0pt}
\setlength{\parskip}{0.5em}

\newcommand{\promptbox}[1]{%
  \begingroup\setlength{\fboxsep}{10pt}\noindent
  \colorbox{blue!8}{\parbox{\dimexpr\linewidth-2\fboxsep\relax}{\textbf{Prompt. }#1}}
  \endgroup
}
\newcommand{\resultbox}[1]{%
  \begingroup\setlength{\fboxsep}{10pt}\noindent
  \colorbox{purple!8}{\parbox{\dimexpr\linewidth-2\fboxsep\relax}{\textbf{Result. }#1}}
  \endgroup
}
\newcommand{\examplebox}[1]{%
  \begingroup\setlength{\fboxsep}{10pt}\noindent
  \colorbox{green!8}{\parbox{\dimexpr\linewidth-2\fboxsep\relax}{\textbf{Example. }#1}}
  \endgroup
}

\title{\textbf{Book 1 --- Lecture 1 --- Exercise 2}\\[2mm]
\large Classifying Six-State Laws}
\author{}
\date{}

\begin{document}
\maketitle

\promptbox{Can you think of a general way to classify the laws that are possible
for a six-state system?}

\section*{Setup}
The six states are the faces of a die. A deterministic law maps each state to
exactly one successor.

\section*{Reversible Case (as in the examples)}
In the four textbook examples, each state has exactly one successor and one
predecessor. So the law is a permutation of six labels. Such laws are
classified, up to relabeling, by their cycle type (equivalently, a partition of
6).

There are 11 partitions of 6:
\[
\begin{aligned}
&(6),\ (5,1),\ (4,2),\ (4,1,1),\ (3,3),\ (3,2,1),\\
&(3,1,1,1),\ (2,2,2),\ (2,2,1,1),\ (2,1,1,1,1),\ (1,1,1,1,1,1).
\end{aligned}
\]

\section*{The Four Examples}
\[
\begin{aligned}
1\to2\to3\to4\to5\to6\to1 &\quad\Rightarrow\quad (6),\\
1\to3\to2\to6\to4\to5\to1 &\quad\Rightarrow\quad (6),\\
1\to2\to6\to1,\ \ 3\to4\to5\to3 &\quad\Rightarrow\quad (3,3),\\
1\to1,\ \ 5\leftrightarrow6,\ \ 2\to3\to4\to2 &\quad\Rightarrow\quad (3,2,1).
\end{aligned}
\]

So examples 1 and 2 are in the same class (both a single 6-cycle), despite
different label orderings.

\resultbox{For reversible six-state laws, classification up to relabeling is
exactly the classification by partitions of 6 (cycle types).}

\examplebox{Two distinct-looking update tables can still represent the same
class. For instance, examples 1 and 2 are both one 6-cycle and therefore
equivalent under relabeling.}

\end{document}
